\section{Localisation Accuracy and Practical Value}
\label{sec:localization}

This section evaluates the accuracy of the radio-fingerprint-based localisation
methods and the channel charting dimensionality-reduction analysis, then
interprets those numbers in plain terms to guide practical deployment decisions.

% ─────────────────────────────────────────────────────────────────────────────
\subsection{What the Numbers Mean for a Real User}
\label{sec:accuracy_meaning}

Localisation accuracy is typically reported as a Mean Absolute Error (MAE) or
90th-percentile error (P90).  To make these figures concrete,
Fig.~\ref{fig:loc_scale} maps MAE and P90 values onto everyday objects and
distances: a \emph{parking space} is $\approx 5\,\text{m}$, a \emph{city block}
is $\approx 100\,\text{m}$, and a \emph{corridor width} is $1{-}2\,\text{m}$.

\begin{figure}[htbp]
  \centering
  \includegraphics[width=\columnwidth]{pictures/loc_scale_visual.png}
  \caption{Real-world scale interpretation of localisation MAE and 90th-percentile
           error for the three fingerprinting methods.  Arrows span MAE (left) to
           P90 (right) on a common distance axis annotated with familiar objects.}
  \label{fig:loc_scale}
\end{figure}

The use-case threshold analysis in Fig.~\ref{fig:loc_cdf} overlays four
industry reference points on the cumulative error distribution:
(i)~sub-metre accuracy required for indoor augmented-reality overlays;
(ii)~$\leq 5\,\text{m}$ required for outdoor pedestrian turn-by-turn guidance;
(iii)~$\leq 15\,\text{m}$ mandated by the EU E-112 emergency-call standard; and
(iv)~$\leq 50\,\text{m}$ sufficient for neighbourhood-level network optimisation.

\begin{figure}[htbp]
  \centering
  \includegraphics[width=\columnwidth]{pictures/loc_usecase_cdf.png}
  \caption{Cumulative distribution of localisation error for all three fingerprinting
           methods, annotated with real-world use-case accuracy thresholds.
           The green band ($\leq 5\,\text{m}$) marks the pedestrian-navigation zone;
           the orange band ($5{-}15\,\text{m}$) marks the emergency-call compliance
           zone.}
  \label{fig:loc_cdf}
\end{figure}

% ─────────────────────────────────────────────────────────────────────────────
\subsection{Fingerprint Localisation Results}
\label{sec:fp_results}

The fingerprint database covers a 1.9-hectare outdoor area
on the Otaniemi campus, sampled on a 10-metre grid
($N = 168$ training positions, 72 test positions)
using 9 base stations at 3.6\,GHz.
Each fingerprint combines OFDM channel-magnitude features, time-difference of
arrival (TDoA) features across 9 BSes, and angle-of-arrival (AoA)
features --- 7191 dimensions in total.

Table~\ref{tab:fp_results} summarises the localization performance.
Fig.~\ref{fig:fp_comparison} shows MAE, RMSE, median and 90th-percentile error
side by side.

\begin{table}[htbp]
  \centering
  \caption{Fingerprint localisation performance on the Otaniemi 100-m scene.}
  \label{tab:fp_results}
  \begin{tabular}{lcccc}
    \toprule
    Method & MAE (m) & RMSE (m) & Median (m) & P90 (m) \\
    \midrule
    WKNN (IDW), $k=2$ & 28.17 & 38.31 & 20.15 & 55.86 \\
    NN Regression     & 28.31  & 32.26  & 25.93  & 52.88  \\
    NN Classification & 74.10  & 85.76  & 71.06  & 132.78  \\
    \bottomrule
  \end{tabular}
\end{table}

\begin{figure}[htbp]
  \centering
  \includegraphics[width=\columnwidth]{pictures/loc_method_comparison.png}
  \caption{Fingerprint localisation accuracy: MAE vs.\ RMSE (left) and median
           vs.\ 90th-percentile error (right) for the three methods on the
           72-sample test set.}
  \label{fig:fp_comparison}
\end{figure}

\paragraph{Plain-language verdict.}
\begin{itemize}
  \item \textbf{WKNN (IDW):} MAE = 28.2\,m — coarse — sufficient for large-area coverage analytics only.
        Provides fast inference with no training cost beyond database construction.
  \item \textbf{NN Regression:} MAE = 28.3\,m — coarse — sufficient for large-area coverage analytics only.
        End-to-end deep network with a continuous 2-D position output; benefits
        from non-linear feature interactions.
  \item \textbf{NN Classification:} MAE = 74.1\,m — coarse — sufficient for large-area coverage analytics only.
        Maps positions to 10-m grid cells; classification accuracy
        0.0\,\% with centroid decoding.
\end{itemize}
The best-performing method (\textbf{WKNN (IDW)}, MAE = 28.17\,m)
meets the E-112 emergency-call requirement of $\leq 50\,\text{m}$ for
85\,\% of test positions, and the pedestrian-navigation
threshold of $\leq 5\,\text{m}$ for 4\,\% of positions.

% ─────────────────────────────────────────────────────────────────────────────
\subsection{Channel Charting}
\label{sec:cc_results}

Channel charting reduces the high-dimensional CSI fingerprint to a
2-D \emph{radio map} without any position labels, using only the
intrinsic geometry of the channel space.  We evaluate 12
method--representation combinations (PCA, MDS, classical MDS, Laplacian
Eigenmaps, Isomap, t-SNE applied to full CSI, inter-BS covariance and
geometric features) and rank them by the Nearest-Neighbour Localisation
Error (NNLE): the mean position error when 5-nearest chart-space neighbours
vote on a location via inverse-distance weighting.

The best channel charting method (t-SNE [RI]) achieves a k-NN localisation error of 23.1\,m, comparable to the supervised fingerprinting baseline — revealing that the radio channel structure encodes geographic position even without labels.

Table~\ref{tab:cc_results} lists all methods.

\begin{table}[htbp]
  \centering
  \caption{Channel charting evaluation summary.}
  \label{tab:cc_results}
  \begin{tabular}{lcccc}
    \toprule
    Method \& TW ($\uparrow$) \& CT ($\uparrow$) \& KS ($\downarrow$) \& NNLE (m) \\
    \midrule
    PCA [RI] \& 0.672 \& 0.837 \& 0.469 \& 42.22 \\
    PCA [Cov] \& 0.590 \& 0.799 \& 0.446 \& 51.74 \\
    PCA [Feat] \& 1.000 \& 1.000 \& 0.000 \& 97.90 \\
    MDS [RI] \& 0.666 \& 0.775 \& 0.259 \& 46.03 \\
    MDS [Cov] \& 0.656 \& 0.794 \& 0.254 \& 55.67 \\
    MDS [Feat] \& 0.035 \& 0.046 \& 0.000 \& 63.30 \\
    CMDS [RI] \& 0.672 \& 0.837 \& 0.469 \& 42.22 \\
    LE [RI] \& 0.593 \& 0.803 \& 1.000 \& 43.10 \\
    Isomap [RI] \& 0.704 \& 0.881 \& 0.424 \& 39.88 \\
    t-SNE [RI] \& 0.816 \& 0.849 \& 0.888 \& 23.14 \\
    t-SNE [Cov] \& 0.728 \& 0.753 \& 0.522 \& 43.17 \\
    t-SNE [Feat] \& 0.648 \& 0.251 \& 11.938 \& 66.87 \\
    \bottomrule
  \end{tabular}
\end{table}


\begin{figure}[htbp]
  \centering
  \includegraphics[width=\columnwidth]{pictures/cc_vs_fingerprinting.png}
  \caption{Channel charting NNLE for all method--representation combinations
           versus the supervised fingerprinting baseline (dashed line).
           Methods below the dashed line match or exceed the unsupervised
           baseline; methods above it still provide useful relative position
           maps without labels.}
  \label{fig:cc_vs_fp}
\end{figure}

\paragraph{Why channel charting matters for operators.}
Even when GPS is unavailable and no labelled database exists --- for example,
during network installation or after a hardware change --- a channel chart can
be constructed from uplink CSI alone.  This is valuable for:
(i)~\emph{initial beam management}: cluster UEs by chart proximity rather than
     explicit location;
(ii)~\emph{anomaly detection}: sudden chart-position shifts indicate abrupt
     propagation changes (new obstacles, antenna failures); and
(iii)~\emph{relative handover prediction}: UE trajectory in chart space
     predicts impending cell-edge crossing without GPS or infrastructure
     modifications.

